\setcounter{section}{2}
\setlength{\parindent}{0pt}  % niente rientro
\setlength{\parskip}{6pt}
\section{Carico e distribuzione del lavoro}

Di seguito viene riportata una stima del carico di lavoro, espressa in ore, sostenuto da ciascun componente del gruppo per la realizzazione dei vari deliverable e del prodotto finale.

\begin{table}[H]
\centering
\renewcommand{\arraystretch}{1.5}
\begin{tabular}{|l|c|c|c|c|c|c|}
\hline
\textbf{Componente} & \textbf{D1} & \textbf{D2} & \textbf{D3} & \textbf{D4} & \textbf{Sviluppo WebApp} & \textbf{TOTALE} \\ \hline
Enrico Sartori & 25 & 20 & 10 & 8 & 300 & \textbf{363} \\ \hline
Stefano Pezzetta & 45 & 5 & 55 & 12 & 80 & \textbf{197} \\ \hline
\textbf{Totale Ore} & \textbf{70} & \textbf{25} & \textbf{65} & \textbf{20} & \textbf{380} & \textbf{560} \\ \hline
\end{tabular}
\caption{Distribuzione del carico di lavoro}
\label{tab:ore_lavoro}
\end{table}

\subsection*{Analisi della distribuzione}
\textbf{Nota:} il monteore indicato indica una stima. Non operando con regimi di orario fissi risulta difficile calcolare il numero esatto di ore dedicato ad ogni sezione.

E' importante sottolineare che le colonne della tabella non si indicano solamente il tempo impiegato nella stesura dei documenti, ma includono l'intero sforzo analitico e progettuale correlato a quelle specifiche aree tematiche. Nello specifico:

\begin{itemize}
    \item \textbf{D1:} comprende anche la creazione del pitch, l'analisi del dominio e la definizione del perimetro del progetto.
    \item \textbf{D3:} include tutte le ore dedicate all'effettiva modellazione architetturale, logica del sistema.
    \item \textbf{D4:} accorpa anche il tempo impiegato per la stesura del Business Plan finanziario e per la realizzazione (registrazione e montaggio) del video multimediale di presentazione.
\end{itemize}

\textbf{Nel dettaglio}

\begin{itemize}
    \item \textbf{Fase D1 (Descrizione ed analisi dei requisiti):} Questa fase ha richiesto un impegno congiunto significativo. La definizione dei requisiti e degli scenari d'uso è stata frutto di numerosi perfezionamenti, necessari per calibrare gli obiettivi di progetto in funzione delle scadenze temporali e delle risorse disponibili, garantendo così la fattibilità della proposta.
    
    \item \textbf{Fase D2 (Sviluppo e Tecnologie):} In questo frangente si evidenzia un carico di lavoro prevalente per \textbf{Enrico Sartori}. Data la natura verticalmente tecnica del documento, la stesura è stata curata quasi interamente da lui per garantire coerenza architetturale e ottimizzare i tempi di produzione.

    \item \textbf{Fase D3 (Analisi e progettazione):} Speculare alla precedente, questa fase è stata guidata principalmente da \textbf{Stefano Pezzetta}, che ha assunto la responsabilità operativa della modellazione UML, della progettazione dello schema concettuale e logico del database, nonché della redazione delle relative specifiche funzionali.

    \item \textbf{Sviluppo WebApp:} L'implementazione del codice ha rappresentato l'attività più onerosa dell'intero ciclo di vita del progetto. Il carico maggiore è stato sostenuto da \textbf{Enrico Sartori}, che ha governato la complessità dello sviluppo Full Stack e l'integrazione delle API, aiutato da \textbf{Stefano Pezzetta} per la realizzazione della componentistica Frontend e della UI.

    \item \textbf{Fase D4 (Reportistica e materiale di presentazione):} La fase conclusiva ha visto un riallineamento dell'impegno lavorativo, con una suddivisione equilibrata delle attività legate alla produzione dei contenuti di presentazione finale, alla stesura del Business Plan e del report finale.
\end{itemize}

